\par\Needspace{10\baselineskip}\pdfbookmark[1]{Prasse. 1814 December 19.}{prasse1814}
\begin{center}
\rule{0.63\linewidth}{0.3pt}\\
Göttingische gelehrte Anzeigen. 1814 December 19.\\
\rule{0.63\linewidth}{0.3pt}
\end{center}

\emph{Tables logarithmiques pour les nombres, les sinus et les tangentes, disposées
dans un nouvel ordre par M. \textup{\textsc{de Prasse}}, professeur des mathématiques à Berlin} (zu
Leipzig) \emph{corrigées et précédées d’une introduction traduite de l’allemand et accompagnée de notes et d’un avertissement par} M. \textsc{Halma} 1814. Paris. 80 Seiten. Preis
Ein Frank.

Ein neuer Abdruck der geschmeidigen von \textsc{Prasse}’schen Tafeln, welche wir
in diesen Blättern [1811 Mai 25] angezeigt haben. Unser dortiges Urtheil über
die von dem französischen Herausgeber unverändert beibehaltene Anordnung der
Tafeln haben wir durch einen drei Jahre länger fortgesetzten Gebrauch derselben
in allen Stücken bestätigt gefunden. In Ansehung der Schönheit des Drucks und
Papiers scheint uns die Französische Ausgabe der Deutschen eher nachzustehen;
doch sind mehrere Druckfehler der letztern hier berichtigt. Wenn man übrigens
bei einem Werke dieser Art, das der verstorbene von \textsc{Prasse} gewiss nicht Gewinnes halber, sondern zum Dienste der Wissenschaft auf seine Kosten unternahm,
auch nicht weiter untersuchen will, in wie fern Hr. \textsc{Halma} zu einem neuen Abdrucke berechtigt war, so kann man doch nicht umhin, sich zu wundern, dass
derselbe, aus Besorgniss seinerseits wieder nachgedruckt zu werden, die einzelnen Exemplare mit seinem Namenszuge bezeichnet hat, und die Nachdrucker gerichtlich zu belangen droht.
\SourceLeaf{253}
\par\Needspace{10\baselineskip}\pdfbookmark[1]{Tafel zur bequemern Berechnung. 1812.}{gauss1812log}
\begin{center}
\rule{0.66\linewidth}{0.3pt}\\
Monatliche Correspondenz, herausg. vom Freih. v. \textsc{Zach}. 1812. Nov.\\
\rule{0.66\linewidth}{0.3pt}
\end{center}

\emph{Tafel zur bequemern Berechnung des Logarithmen der Summe oder Differenz
zweier Grössen, welche selbst nur durch ihre Logarithmen gegeben sind. Von Herrn
Prof.} \textsc{Gauss}.

Je weiter sich beständig die Geschäfte der rechnenden Astronomen ausdehnen, desto wichtiger wird ihnen jede, wenn auch an sich nur kleine Erleichterung derselben. Die \emph{Monatliche Correspondenz} hat sich hierin schon vielfältige
Verdienste erworben, indem sie mancherlei Tafeln aufgenommen hat, deren kleiner Umfang nicht verstattete, sie besonders herauszugeben. Ich lege daher gern
in derselben eine kleine Tafel nieder, die freilich nicht eigentlich astronomisch
ist, aber besonders doch den rechnenden Astronomen willkommen sein wird, und
die etwa in Zukunft sehr zweckmässig mit einem neuen Abdruck der kleinen \textsc{La
Lande}’schen Tafeln verbunden werden könnte. Das Geschäft, was sie erleichtern soll, kommt bei astronomischen Rechnungen alle Augenblick vor; es erfordert sonst ein dreimaliges, oder wenn man eine leichte Verwandlung anwendet,
doch nothwendig ein zweimaliges Aufschlagen in den Logarithmen-Tafeln, was
hier auf ein einziges gebracht wird. Die Idee dazu hat \textsc{Leonelli}, so viel ich weiss,
zuerst angegeben; allein seine Meinung war, eine solche Tafel für Rechnungen
mit 14 Decimalen zu construiren, und gerade dies kann ich nicht zweckmässig
finden. Sie würde bei einer solchen Ausdehnung einen grossen Folioband füllen,
ihre Berechnung würde eine ungeheuere Arbeit und Zeit erfordern, und sie würde
fast nie von Nutzen, und immer nur von wenig Nutzen sein, da so scharfe Rechnungen so selten --- in der eigentlichen practischen Astronomie nie --- vorkommen, dass die verhältnissmässig doch nur kleine Erleichterung die Construction,
ja nicht einmal den Ankauf einer solchen Tafel belohnen würde. Ich habe diese
Tafel zu meinem eigenen Gebrauch für Rechnungen mit 5 Decimalen, die in der
Ausübung die häufigsten sind, schon vor vielen Jahren construirt, und die, wenn
auch \emph{jedesmal} kleine, doch wenn sie viele Tausendmale wiederkehrt, sehr erhebliche Erleichterung, hat mir die darauf gewandte Mühe bereits reichlich ersetzt.
Es wäre zu wünschen, dass jemand sich der Arbeit unterzöge, eine ähnliche Tafel in 10 oder 100 mal so grosser Ausdehnung für Rechnungen mit 7 Decimalen zu
\SourceLeaf{254}
construiren, die als ein sehr schätzbares Supplement den gewöhnlichen Logarithmen-Tafeln beigefügt werden könnte.

Die Einrichtung der aus drei Columnen bestehenden Tafel ist sehr einfach.
Die erste Columne \(A\) geht von 0 bis 2 durch alle Tausendtheile, von da bis \(3,4\)
durch alle Hunderttheile, und von \(3,4\) bis \(5,0\) durch alle Zehntheile; mit \(5,0\)
kann die Tafel für 5 Decimalen als geschlossen angesehen werden, da die zweite
Columne für diesen und für grössere Werthe von \(A\) verschwindet, und die Zahlen der dritten Columne denen der ersten gleich werden. Setzt man eine Zahl
der ersten Columne \(A=\log m\), so ist in der zweiten Columne \(B=\log(1+\frac{1}{m})\)
und in der dritten Columne \(C=\log(1+m)\), so dass immer \(C=A+B\). Man
kann also auch die Zahlen der drei Columnen als die doppelten Logarithmen der
Tangenten, Cosecanten und Secanten der Winkel von \(45^{0}\) bis \(90^{0}\) betrachten.
Die Anwendung davon ist nun folgende:

\begin{center}
I.\quad \emph{Aus den Logarithmen zweier Grössen \(a,b\) den Logarithmen der Summe
zu finden.}
\end{center}

Es sei \(\log a\) der grössere Logarithm., man gehe mit \(\log a-\log b\) in die
Columne \(A\) ein, und nehme daneben entweder aus der zweiten Columne \(B\), oder
aus der dritten Columne \(C\). Man hat dann
\[
\begin{aligned}
&\log(a+b)=\log a+B\\
\text{oder}\qquad &\log(a+b)=\log b+C.
\end{aligned}
\]

\begin{center}
II.\quad \emph{Aus den Logarithmen zweier Grössen \(a,b\) den Logarithmen der Differenz
zu finden.}
\end{center}

\emph{Erstens,} ist die Differenz der Logarithmen \(\log a-\log b\) grösser als \(0,30103\),
so suche man dieselbe in \(C\), wodurch man hat
\[
\begin{aligned}
&\log(a-b)=\log a-B\\
\text{oder}\qquad &\log(a-b)=\log b+A
\end{aligned}
\]

\emph{Zweitens,} ist \(\log a-\log b\) kleiner als \(0,30103\), so suche man es in \(B\), wodurch wird
\[
\begin{aligned}
&\log(a-b)=\log a-C\\
\text{oder}\qquad &\log(a-b)=\log b-A
\end{aligned}
\]
\SourceLeaf{255}
Es gibt daher bei jeder Aufgabe zwei Auflösungsarten; man thut aber wohl,
sich an eine bestimmte zu gewöhnen, um sich den Gebrauch der Tafel desto leichter mechanisch zu machen. Mir ist dies bei der jedesmal zuerst angesetzten Manier am bequemsten gefallen.

\Needspace{12\baselineskip}\begin{center}\textit{Beispiele:}\end{center}

I. Aus $\log a=0{,}36173$ und $\log b=0{,}23045$ den Logarithmen der Summe
zu finden, sucht man $0{,}13128$ in $A$, wobei man findet
\[
\begin{array}{rcr}
B&\cdot\;\cdot\;\cdot\;\cdot\;\cdot&0{,}24033\\
\log a&\cdot\;\cdot\;\cdot\;\cdot\;\cdot&0{,}36173\\\hline
\log(a+b)&\cdot\;\cdot&0{,}60206
\end{array}
\qquad
\begin{array}{rcr}
C&\cdot\;\cdot\;\cdot\;\cdot\;\cdot&0{,}37161\\
\log b&\cdot\;\cdot\;\cdot\;\cdot\;\cdot&0{,}23045\\\hline
&&0{,}60206
\end{array}
\]

II. Aus $\log a=0{,}89042$, und $\log b=0{,}24797$ den Logarithm. der Differenz zu finden. Da $\log a-\log b=0{,}64245$ grösser als $0{,}30103$, so sucht
man es in der Columne $C$, woneben man findet
\[
\begin{array}{rcr}
B&\cdot\;\cdot\;\cdot\;\cdot\;\cdot&0{,}11227\\
\log a&\cdot\;\cdot\;\cdot\;\cdot\;\cdot&0{,}89042\\\hline
\log(a-b)&\cdot\;\cdot&0{,}77815
\end{array}
\qquad
\begin{array}{rcr}
A&\cdot\;\cdot\;\cdot\;\cdot\;\cdot&0{,}53018\\
\log b&\cdot\;\cdot\;\cdot\;\cdot\;\cdot&0{,}24797\\\hline
&&0{,}77815
\end{array}
\]

III. Aus $\log a=0{,}25042$, $\log b=0{,}19033$ den Logarithmen der Differenz zu finden. Hier gibt $\log a-\log b=0{,}06009$ in $B$ aufgesucht.
\[
\begin{array}{rcr}
C&\cdot\;\cdot\;\cdot\;\cdot\;\cdot&0{,}88871\\
\log a&\cdot\;\cdot\;\cdot\;\cdot\;\cdot&0{,}25042\\\hline
\log(a-b)&\cdot\;\cdot&9{,}36171
\end{array}
\qquad
\begin{array}{rcr}
A&\cdot\;\cdot\;\cdot\;\cdot\;\cdot&0{,}82862\\
\log b&\cdot\;\cdot\;\cdot\;\cdot\;\cdot&0{,}19033\\\hline
&&9{,}36171
\end{array}
\]

\par\Needspace{10\baselineskip}\pdfbookmark[1]{Pasquich. 1817 October 4.}{pasquich1817}
\begin{center}
\rule{0.65\linewidth}{0.4pt}\\
Göttingische gelehrte Anzeigen. 1817 October 4.\\
\rule{0.65\linewidth}{0.4pt}
\end{center}

\textit{Abgekürzte Logarithmisch-Trigonometrische Tafeln, mit neuen Zusätzen zur
Abkürzung und Erleichterung trigonometrischer Rechnungen, herausgegeben von }\textsc{Joh.}
\textsc{Pasquich}\textit{, Director der Königl. Ofner Sternwarte.} Leipzig. In der Weidmannschen
Buchhandlung 1817. XXXII und 228 Seiten in Octav. (Auch mit Lateinischem Titel.)
\par\SourceLeaf{256}
Kleinere logarithmische Tafeln, mit fünf Decimalen, sind bei denjenigen,
die viel mit Zahlenrechnungen zu verkehren haben, besonders bei den Astronomen, sehr beliebt, weil in der That die Fälle, wo sie ausreichen, häufig. ja die häufigeren, sind, und durch ein bequemes Format und eine mässige Grösse die Arbeit
sehr erleichtert wird. Die kleinen netten \textsc{Lalande}'schen und die \textsc{Prasse}'schen
Tafeln sind in Jedermanns Händen; bei letztern ist das Zusammendrängen in einen kleinen Raum so weit wie möglich getrieben, zum Theil aber allerdings auf
Kosten der Bequemlichkeit. Die Herausgabe der vorliegenden auch nur auf fünf
Stellen gehenden Tafeln ist, wie in der Vorrede berichtet wird, durch den von
\textsc{Gauss} in der monatlichen Correspondenz 1812 geäusserten Wunsch veranlasst,
dass die daselbst zuerst abgedruckte Tafel zur bequemen Berechnung der Logarithmen der Summen und Differenzen einer neuen Ausgabe der \textsc{Lalande}'schen Tafeln einverleibt werden möchte. Der neue Abdruck dieser Hülfstafel in gegenwärtiger Sammlung wird denjenigen angenehm sein, denen der erste Abdruck
nicht zu Gebote stand, oder denen der Gebrauch derselben in der M. C. zu beschwerlich war. Ausserdem zeichnet sich diese Sammlung noch durch eine neue
von Hrn. \textsc{Pasquich} berechnete den trigonometrischen Tafeln beigefügte Hülfstafel aus, deren wir unten mit mehrern erwähnen werden.

Die Logarithmen der Zahlen gehen, wie bei \textsc{Lalande} und von \textsc{Prasse} bis
10000, und sind, so wie bei jenen, hier in ihrer natürlichen Ordnung gedruckt.
Doch vermisst man ungern ein Paar Erleichterungsmittel, welche bei den \textsc{Lalande}'schen Tafeln Statt finden; es sind nemlich theils die \textit{Differenzen} nicht beigefügt, theils die untersten Logarithmen jeder Spalte oben in der nächstfolgenden
nicht wiederholt. In den \textsc{Prasse}'schen Tafeln findet man zwar diese Bequemlichkeit auch nicht, allein dort werden sie durch den kleinen Raum der Tafel mehr
als ersetzt, da jene auf 24 Seiten eben dasselbe liefern, was bei \textsc{Pasquich} auf 56
Seiten eines beträchtlich grössern Formats steht. Dies scheinen zwar nur Kleinigkeiten, und sie sind es auch für alle, die nur dann und wann einmal Logarithmen aufzuschlagen haben, aber nicht für solche, die Logarithmen-Tafeln beständig zur Hand haben müssen.

In den trigonometrischen Tafeln enthält immer jede Seite zur linken die
Logarithmen der Sinus, Cosinus, Tangenten und Cotangenten, und zwar so,
dass je drei Seiten zwei Grade fassen. Diese Einrichtung, welche durch das gewählte Format und die Schrift herbeigeführt wurde, scheint uns etwas unbequem;
\SourceLeaf{257}
wir hätten entweder ein kleineres Format, immer mit einem halben Grad auf der
Seite, oder ein etwas weniges längeres mit kleinerer Schrift, so dass ein ganzer
Grad auf die Seite gekommen wäre (wie in \textsc{Shervins} Tafeln) vorgezogen. Von diesen Logarithmen sind immer nur die vier, drei oder zwei letzten Ziffern, so lange
die vorhergehenden ungeändert bleiben, abgedruckt, wodurch dem Copiisten, dem
Setzer und dem Corrector die Arbeit erleichtert wurde, und die Tafeln ein reinlicheres Ansehen erhalten: dem ungeachtet können wir diese Einrichtung bei
Tafeln, die zum täglichen Gebrauch bestimmt sind, nicht unbedingt billigen, da
das Auge immer die, wenn auch nur kleine, Mühe hat, in der Columne erst in
die Höhe zu gehen, und die übrigen Ziffern zu finden. Die Differenzen der Logarithmen findet man hier sogleich mit 60 dividirt; eine Einrichtung, welche auch
in einigen andern Tafeln gewählt ist, in der Absicht, das Interpoliren zu erleichtern. Ob diese Erleichterung wirklich Statt findet, oder nicht, wird von der Gewöhnung des Rechners abhängen. Rec. findet in dieser Beziehung die \textsc{Lalande}schen Tafeln, wo die ganzen Differenzen angesetzt sind, wenigstens nicht unbequemer. Bei der Kleinheit der Zahlen, mit denen zu operiren ist, macht ein
etwas geübter Rechner die zum Behuf des Interpolirens nöthigen Operationen
leicht im Kopfe, und findet fast immer diesen oder jenen Local-Vortheil zu benutzen Gelegenheit. Dabei hat man noch die angenehme Gewissheit, sein Interpolations-Resultat so scharf zu erhalten, als es möglich ist; bei der von Hrn. \textsc{Pasquich} gewählten Einrichtung hingegen ist, allgemein zu reden, der Fehler des
Interpolirens etwas grösser, welches indessen ausführlicher zu entwickeln hier
nicht der Ort ist.

Die Seite zur rechten enthält bei den trigonometrischen Tafeln die Quadrate
der Sinus, Cosinus, Tangenten und Cotangenten, welche zur Erleichterung des
Interpolirens dienen sollen, wenn man aus dem Logarithmen eines Sinus, Cosinus, einer Tangente oder Cotangente den Logarithmen einer der drei andern trigonometrischen Functionen verlangt, ohne den Bogen selbst nöthig zu haben.
Diese Operation kömmt allerdings äusserst häufig vor, und das gewöhnliche Verfahren erfordert beim Interpoliren eine Multiplication und eine Division, wo mit
Hrn. \textsc{Pasquich}'s Hülfstafel eine Multiplication ausreicht. Es ist nemlich, für das
Interpoliren hinreichend genau,
\SourceLeaf{258}
\[
\begin{aligned}
\Delta\log\cos\varphi&=-\operatorname{tang}\varphi^{2}.\Delta\log\sin\varphi\\
\Delta\log\operatorname{tang}\varphi&=-\Delta\log\operatorname{cotang}\varphi=(1+\operatorname{tang}\varphi^{2}).\Delta\log\sin\varphi\\
\Delta\log\sin\varphi&=-\operatorname{cotang}\varphi^{2}.\Delta\log\cos\varphi\\
\Delta\log\operatorname{tang}\varphi&=-\Delta\log\operatorname{cotang}\varphi=-(1+\operatorname{cotang}\varphi^{2}).\Delta\log\cos\varphi\\
\Delta\log\sin\varphi&=\cos\varphi^{2}.\Delta\log\operatorname{tang}\varphi=-\cos\varphi^{2}.\Delta\log\operatorname{cotang}\varphi\\
\Delta\log\cos\varphi&=-\sin\varphi^{2}.\Delta\log\operatorname{tang}\varphi=\sin\varphi^{2}.\Delta\log\operatorname{cotang}\varphi
\end{aligned}
\]
Inzwischen muss Rec. gestehen, dass er dem ungeachtet das gewöhnliche Verfahren zum Interpoliren nicht bloss eben so bequem, sondern sogar bequemer findet. Theils wird es immer erst einige Mühe kosten, sich die obigen sechs Formeln so mechanisch zu machen, dass man sie ohne alles Besinnen oder ohne ein
besonderes Blatt neben sich zu legen, richtig anwendet; theils ist es beschwerlich, den Multiplications-Factor erst auf der andern Seite aufzusuchen, oder vielmehr zusammen zu suchen, da die oben erwähnte Trennung der ersten und letzten Ziffern auch hier beim Abdruck gewählt ist; endlich hat man bei dem gewöhnlichen Verfahren es immer nur mit kleinen Zahlen zu thun, mit denen man
leicht im Kopf rechnet, da hingegen die Quadrate in \textsc{Pasquich}'s Tafeln mit fünf
Decimalen angesetzt sind, die man freilich nicht alle braucht, aber die gerade deswegen, wie jeder erfahrne Rechner weiss, störend sind. Ausserdem können wir
hier nicht unerwähnt lassen, dass das gewöhnliche Verfahren, allgemein zu reden,
\textit{schärfer} ist, als diese künstlichere Interpolation (die Gründe dieser Behauptung,
von der man vielleicht bei einer weniger genauen Prüfung gerade das Gegentheil
glauben könnte, würden für diesen Ort zu weitläuftig sein). Wir begnügen uns
das Gesagte bloss durch ein Beispiel zu erläutern. Soll zu \(\log\cos\varphi=9{,}92478\)
der \(\log\operatorname{tang}\varphi\) gesucht werden, so findet man den Proportionaltheil aus \textsc{Pasquich}'s Tafel durch die Berechnung von \(4\times(1+2{,}4170)=13{,}668\) oder am nächsten \(=14\), also \(\log\operatorname{tang}\varphi=9{,}80850\), während die gewöhnliche Methode den
Proportionaltheil eben so bequem durch die Entwicklung von \(\frac{4\times28}{9}=12\frac49\), am
nächsten \(=12\), und den gesuchten Logarithmen \(=9{,}80848\) gibt. In diesem
Beispiele ist auch das Resultat der gewöhnlichen Methode das schärfere; in andern Fällen kann auch das umgekehrte Verhältniss Statt finden, aber im Durchschnitt wird der Vortheil in dieser Beziehung auf Seiten des gewöhnlichen Verfahrens sein. Uebrigens wollen wir nicht in Abrede stellen, dass dieser Theil
der Tafel, wenn auch das Interpoliren nicht dadurch gewinnt, doch zuweilen für
\SourceLeaf{259}
andere Zwecke angenehm sein könne; allein die Bequemlichkeit logarithmischer
Handtafeln, die man zum täglichen Gebrauch bestimmt, verliert natürlich in demselben Verhältniss, als ihr Umfang vergrössert wird. Wir bemerken noch, dass
in dem ersten Grade die trigonometrischen Logarithmen von \(10\) zu \(10\) Secunden
bis \(56\) Minuten, und in den vier letzten Minuten von \(20\) zu \(20\) Secunden angesetzt sind.

Die \textsc{Gauss}ische Tafel für die Logarithmen der Summen und Differenzen ist
ganz unverändert abgedruckt. Inconsequent scheint es uns aber zu sein, wenn
der Verf. in der Einleitung den Nutzen einer ähnlichen Tafel mit sieben Decimalen in Zweifel zieht. Ist anders eine solche Tafel zweckmässig eingerichtet, so ist
ihr Nutzen bei scharfen Rechnungen gerade eben so gross, als der Nutzen der
hier wieder abgedruckten Tafeln bei Rechnungen mit fünf Decimalen: bei den
kleinern Tafeln, eben so wie bei den grössern, wird der dadurch zu erhaltende
Zeitgewinn natürlich nur solchen Personen fühlbar, die \textit{viel} zu rechnen haben.
Wir haben jetzt bald die Erscheinung einer solchen grössern Tafel, von einer geschickten Hand berechnet, zu erwarten.

\par\Needspace{10\baselineskip}\pdfbookmark[1]{Matthiessen. 1819 Januar 30.}{matthiessen1819}\medskip\hrule\smallskip
\begin{center}
Göttingische gelehrte Anzeigen. 1819 Januar 30.
\end{center}
\smallskip\hrule\medskip

[E. A. \textsc{Matthiessen}.] \textit{Tafel zur bequemern Berechnung des Logarithmen der
Summe oder Differenz zweier Grössen, welche selbst nur durch ihre Logarithmen gegeben sind.} Altona 1818. Bei J. F. Hammerich. Einleitung 33 S. Die Tafeln
212 S. in Quart. (Titel und Einleitung auch in lateinischer Sprache.)

Bei etwas ausgedehnten Rechnungen ist jede Abkürzung schätzbar; auch
solche Hülfsmittel, die bei einer einmaligen Anwendung nur einen kleinen Vortheil gewähren, werden durch oft wiederkehrende Benutzung wichtig. Ein solches Erleichterungsmittel ist eine im Jahr 1812 in der monatlichen Correspondenz
zuerst gegebene Hülfstafel, um aus den Logarithmen zweier Grössen unmittelbar
die Logarithmen ihrer Summe oder Differenz abzuleiten: man erreicht dadurch
mit Einem Aufschlagen, wozu man sonst ein dreimaliges oder wenigstens zweimaliges nöthig hätte. Im Besitz einer solchen Tafel wird man mit Vortheil manche
\SourceLeaf{260}
Formeln in ihrer ursprünglichen Gestalt beibehalten können, denen man sonst
wohl, durch Einführung von Hülfswinkeln, eine zur Rechnung bequemere Form
zu geben sucht. Die erwähnte Tafel war nur für Rechnungen mit fünf Decimalen
bestimmt: eben weil solche Rechnungen am häufigsten vorkommen, lag dies Bedürfniss am nächsten. Der bei Bekanntmachung derselben geäusserte Wunsch,
dass jemand sich der Mühe unterziehen möchte, eine ähnliche Tafel in grösserm
Umfange und mit sieben Decimalen zu berechnen, hat die vorliegenden Tafeln
veranlasst, deren Verfasser, Hr. \textsc{Matthiessen}, sich durch diese mühsame Arbeit
ein Recht auf den Dank aller derer erworben hat, die viel mit logarithmischen
Rechnungen zu thun haben. Die Vorrede gibt Nachricht von der Methode, deren sich der Verfasser zur Berechnung bedient hat, und von seiner lobenswerthen
Sorgfalt, die Tafel auch in der letzten Ziffer durchgehends zuverlässig zu machen.
Bei Tafeln dieser Art ist auch die äussere Einrichtung keinesweges gleichgültig.
Der Verf. hatte anfangs eine Anordnung im Sinn, deren Zweck war, die Tafel in
den möglich kleinsten Raum zusammenzudrängen. Allein da die ganze Bestimmung der Tafel nur dahin geht, die Rechnungen zu \textit{erleichtern}, so würde jene
ganz verfehlt werden, wenn die Anordnung der Tafel zu künstlich wäre, und eine
beschwerliche Aufmerksamkeit erforderte. Der Verf. entschloss sich daher bei
reiferer Ueberlegung mit Recht, jene wenn gleich sinnreiche Anordnung bei Seite
zu setzen, und statt derselben eine einfachere zu wählen, obgleich der Umfang
des Bandes dadurch beträchtlich vergrössert wurde. Vielleicht wird mancher,
der die Tafeln gebraucht, mit uns wünschen, dass der Verf. hierin lieber noch
etwas weiter gegangen wäre, und die Columne \(B\) und \(C\) jede vollständig hätte
abdrucken lassen, deren vier letzte, beiden gemeinschaftliche, Ziffern nur einmal
dastehen. Der Ueberblick würde dadurch noch bequemer geworden sein, und
das Format wäre auch dadurch nicht vergrössert, wenn etwas kleinere Schrift gewählt wäre, welches ohne Nachtheil, vielleicht selbst mit Vortheil für das gefällige Ansehen, hätte geschehen können. Auch die fast zu strenge Oekonomie mit
den Ziffern, wo in der Regel nur die vierte abgedruckt ist, so lange die vorhergehenden ungeändert bleiben, und mit den Proportionaltheilen, die immer nur
Einmal angesetzt sind, und also zuweilen ein Zurück- oder Vorausblättern nöthig
machen, thut der Bequemlichkeit einigen Eintrag. Endlich hätten wir gewünscht,
dass der letzte Theil der Tafel S. 212 in einer zehnmal. grössern Ausdehnung gegeben wäre, um die zweiten Differenzen unmerklich zu machen; es wäre dazu
\SourceLeaf{261}
nur Eine Seite mehr erforderlich gewesen. Alle diese Bemerkungen, die zum
Theil mit auf individueller Gewöhnung beruhen mögen, sollen das Verdienstliche
dieser Arbeit keinesweges schmälern, welches gewiss von allen anerkannt wird,
die von derselben Gebrauch zu machen Gelegenheit nehmen werden.

\par\Needspace{12\baselineskip}\pdfbookmark[1]{Gauss an Schumacher. 1822 November 25.}{schumacher1822}\begin{center}\rule{0.72\linewidth}{0.3pt}\par Astronomische Nachrichten, herausg. v. \textsc{Schumacher}. Nr. 24. Beilage 2. 1822 Dec.\par\rule{0.72\linewidth}{0.3pt}\end{center}

\begin{center}
\textsc{Gauss} an \textsc{Schumacher}.
\end{center}
\begin{flushright}
Göttingen 1822. Nov. 25.
\end{flushright}
--- --- Den von Herrn Professor \textsc{Encke} geäusserten Wunsch Logarithmentafeln mit
6 Ziffern betreffend, habe ich schon öfters gehegt, und Sie werden sich gewiss
vielfältigen Dank erwerben, wenn Sie solche veranlassten. Alles, was Herr \textsc{Encke}
über das Aeussere und Innere sagt, unterschreibe ich als meine eigene Meinung,
nur die Proportionaltheile scheinen mir überflüssig, und alles Ueberflüssige schadet dem leichten übersichtlichen Gebrauch. Bei solchen Dingen hängt freilich
manches von individueller Gewöhnung ab; indessen wenn einige, die Gebrauch
von Tafeln machen, anders gewöhnt sind als ich, so sind doch auch wohl andere
eben so gewöhnt; und daher berühre ich noch einen Umstand, nämlich die Abänderung der \(4^{\mathrm{ten}}\) Ziffer für die Logarithmentafeln. Sie kennen die Einrichtung,
die in dieser Beziehung in \textsc{Callet}'s Tafeln gemacht ist, und einige haben dies als
eine Verbesserung betrachtet. Ich gestehe, dass ich der entgegengesetzten Meinung bin, und die \textit{regelmässige} Abtheilung von 5 zu 5 Zeilen durch horizontale
Striche, wie sie in \textsc{Sherwin}'s und andern Tafeln ist, für etwas, bei häufigem Gebrauche \textit{viel} wesentlicheres und bequemeres halte, daher ich mich der \textsc{Callet}'schen Logarithmen auch niemals bedienen mag. Bei meiner vieljährigen Praxis
weiss ich auch nicht einen einzigen Fall, wo der Gebrauch der \textsc{Sherwin}'schen Tafeln mich bei der \(4^{\mathrm{ten}}\) Ziffer zu einem Rechnungsfehler verleitet hätte, daher ich
auch auf die Sternchen bei \textsc{Vega} und andern Tafeln gar keinen Werth lege, und
des bessern Papiers und der schönern Ziffern wegen, mich lieber an die \textsc{Sherwin}'schen halte. Wer auf solche Warnungszeichen einen Werth setzt, kann sich
leicht in seinem Exemplare an den betreffenden Stellen rothe oder grüne Punkte
machen. ---
